\section{Localized modes and self-adjoint eigenproblems}
\label{sec:localized}

A core goal of the substrate theory is to explain particle-like behavior (localization, discrete spectra,
internal ``spin/charge''-like structure) as properties of confined wave states rather than point particles.
In this framework, particles are extended standing-wave configurations of the substrate.

\subsection{Linearization about symmetric backgrounds and self-adjointness}

Let $\vecX_0(x)$ be a time-independent background embedding representing a localized deformation region.
We are especially interested in backgrounds that are \emph{spherically symmetric} in the intrinsic brane coordinates,
so that invariants depend only on $r=|x|$.

Linearize the equations of motion \eqref{eq:eom} about $\vecX_0$ by writing
$\vecX(x,t)=\vecX_0(x)+\bm{\eta}(x,t)$ with small perturbation $\bm{\eta}$.
Because the dynamics derives from an action \eqref{eq:action}, the linearized operator arises from the second
variation of the action and is therefore self-adjoint with respect to the inner product induced by the kinetic energy,
schematically
\begin{equation}
  \langle \bm{\eta},\bm{\zeta}\rangle
  :=\int_{\Omega} \rho_m\,\bm{\eta}(x)\cdot \bm{\zeta}(x)\,\dd^3x,
  \label{eq:inner-product}
\end{equation}
together with the boundary/regularity conditions appropriate to the chosen domain.
Orthogonality and completeness of eigenmodes are therefore not added by hand; they are standard properties of a
self-adjoint eigenproblem.

\subsection{Angular structure and spherical harmonics}

For a spherically symmetric background, the linearized operator commutes with the action of $SO(3)$ on the intrinsic
coordinates.
Consequently, eigenmodes can be chosen to transform in irreducible $SO(3)$ representations.
Operationally, this implies a separation into a radial problem and an angular problem whose eigenfunctions are
spherical harmonics $Y_{\ell m}(\theta,\varphi)$.
This statement is purely a symmetry consequence: it requires rotational invariance of the background and of the
linearized operator, not a change in the underlying substrate ontology.

Appropriate superpositions, together with nontrivial embedding-space polarization structure, can yield torus-like,
loop-like, or shell-like energy density profiles.
These should be understood as standing-wave organizations of the substrate rather than as literal point particles
moving on thin trajectories.

\subsection{Amplitude scaling and energy normalization}

The linear eigenproblem fixes mode shapes and frequencies but not the overall amplitude:
if $\bm{\eta}$ is an eigenmode, so is $c\,\bm{\eta}$ for any scalar $c$.
Physical amplitudes are selected only when nonlinearities are included and the total energy is fixed.
For an electron-like excitation one may impose a rest-energy normalization
\begin{equation}
  \overline{E}[\vecX] \;=\; m_e c^2,
  \label{eq:rest-energy-normalization}
\end{equation}
where $\overline{E}$ denotes the conserved substrate energy associated with \eqref{eq:action}.
For baryon-like modes the same logic applies with the appropriate target rest energy of the chosen species.
The important structural point is unchanged: amplitude is not determined by the linear eigenproblem alone.

\subsection{Confinement as geometric self-guidance}

A confining mechanism is required: in a purely linear wave system, localized packets generally disperse or radiate.
A useful conceptual benchmark is the toroidal electron proposal attributed to Williamson \& van der Mark, which argues
that ordinary gravitational curvature is far too weak to confine a Compton-scale photon-like mode and that confinement
must arise from a self-consistent, excitation-specific mechanism.
We adopt the \emph{confinement requirement} while not committing to a literal double-loop trajectory
(see e.g.\ the review discussion in \citep{HuygensOptics2025}).

In the present substrate theory, the first candidate confinement mechanism is the brane's \emph{geometric}
self-guidance.
Large-amplitude deformations modify the local Euclidean link lengths on the lattice and the induced metric in the
continuum; this changes the local stress response and therefore the propagation characteristics of fluctuations around
the deformed state.
The crucial point is that this mechanism is already present even when the microscopic link law is Hooke-linear in the
scalar extension.
Constitutive refinements may strengthen or weaken localization, but the underlying nonlinear feedback comes from the
same geometry that couples transverse amplitude to lateral strain.

\subsection{Why the baryon sector is the next numerical target}

The cubic substrate naturally singles out three axis-associated channels.
That makes the proton/neutron sector a particularly informative next target: a confined three-component standing-wave
mode tests the very structure that the lattice introduces most directly.
In contrast, an isolated electron-like ansatz can exist without decisively probing the triplet mixing sector.

The working hypothesis is therefore that a baryon-like seed should be built from three mutually locked narrowband
components tied primarily to the axis-aligned carrier channels, with confinement arising from geometric self-guidance
and with the far-field Abelian trace sector distinguishing charged and neutral configurations.
This is not yet a derivation of QCD or of nucleon phenomenology.
It is a concrete question for simulation: does the substrate support stable, localized, color-confined triplet modes at
all?
If the answer is no, an important bridge in the theory fails early and clearly.
If the answer is yes, the cubic triplet sector becomes the natural arena for sharpening the particle interpretation.

\subsection{Connection to Berry/Wilczek--Zee holonomy and spin}

Localized modes can carry internal polarization structure (multi-component standing waves).
The Berry/Wilczek--Zee connections of \Cref{sec:geophase} can then be evaluated on loops surrounding the localized
region or on adiabatic deformations of a localized mode family.
In this way, long-range ``charge''-like labels are interpreted as far-field holonomy/phase structure, while the
spin-$\tfrac12$ $4\pi$ aspect is attributed to spinorial holonomy of a two-level polarization subspace.
